\documentclass[12pt]{article}
\usepackage[utf8]{inputenc}
\usepackage[T1]{fontenc}
\usepackage{hyperref}
\title{North Kodut Language Resources for K-12 Curriculum Development}
\author{Compiled by Claude Sonnet 5 (claude-sonnet-5) for Andr\'as Kornai}
\date{}
\begin{document}
\maketitle
\section{Introduction}
North Kodut (also North Watut, Unank, Onank; ISO 639-3 una) is an Austronesian language of Morobe Province, Papua New Guinea. Speaker estimates vary, from 465 to about 600 in the sources found. Speakers also use Tok Pisin and some English. Very little material exists.
\section{Primary Educational Resources}
No primers, textbooks or school programmes were found in search results.
\section{Linguistic and Reference Materials}
\begin{itemize}
\item Numeral system data for Kodut North (Unank). \url{https://lingweb.eva.mpg.de/channumerals/Kodut-North-Unank.htm}
\item ANU Archives record of SIL survey word lists and grammar notes for North Watut (Unank). \url{https://archivescollection.anu.edu.au/index.php/north-watut-watut-3-language-papua-new-guinea-includes-lexis-grammar-outline-grammar-sketch-syntactic-check-list-notes-sil-survey-word-lists-for-unank-morobe-province-and-maralinan-morobe-province}
\end{itemize}
\section{Religious and Community Resources}
\begin{itemize}
\item Joshua Project page for Onank (Northern Watut). \url{https://joshuaproject.net/index.php/people_groups/14160}
\end{itemize}
\section{Digital and Audio Resources}
\begin{itemize}
\item Global Recordings Network catalogue page for code una. \url{https://globalrecordings.net/en/language/una}
\end{itemize}
The Middle Watut dictionary on Webonary is a different language (mpl) and is not included.
\section{Curriculum Development Context}
No usable curriculum materials were found. Any work would need partnership with the Waffa valley communities, who were not identified in the searches.
\end{document}